\section{Testing}\label{sec:testing}

\subsection{Approach}

The program was tested at two levels.

\paragraph*{Automated tests} JUnit tests~\cite{junit}, run through Spring Boot's test support, start the
application's Spring context without opening a window and call the task service exactly as the screens do, so they
exercise the same rules the user meets. Most test classes extend \texttt{DatabaseTestBase}, which gives each class a
new, empty database in a temporary folder: Flyway runs the migration on it and Hibernate checks the entity against
the table, as when the program starts for the first time. The context is closed after each class, so that its
connection lets go of the file before JUnit deletes the folder, which Windows would otherwise refuse. Two classes
start further contexts of their own: \texttt{RestartTest} starts the program's Spring side on a folder, adds a task,
closes it and starts it again on the same folder, and \texttt{TwoCopiesTest} opens a second context on the same
database to play the second copy of the program in test case 6.

A seventh class, \texttt{TaskTextTest}, needs neither Spring nor a window: it checks the words the screen puts
together, such as how a time is written and what an empty list says, by calling \texttt{TaskText} directly. The tests
are organised in these seven classes:
\begin{itemize}
  \item \texttt{BriefTestCasesTest}: the brief's seven test cases, one test each, in the brief's order;
  \item \texttt{TaskRulesTest}: the rules a task is saved under, beyond the brief's cases;
  \item \texttt{SearchAndCountsTest}: the search, the status filter and the counts;
  \item \texttt{TwoCopiesTest}: two copies of the program on one database;
  \item \texttt{RestartTest}: tasks surviving a restart;
  \item \texttt{FoundationTest}: the program starting on an empty database;
  \item \texttt{TaskTextTest}: the words the window shows.
\end{itemize}
They run with \texttt{./mvnw test}, and \texttt{./mvnw -q test} was run before every commit.

\paragraph*{Tests in the running program} Every test case of the brief was also carried out by hand in the running
program, and the screenshots of Section~\ref{sec:functionality} record what it showed. Test case 6 needs two copies of
the program open on the same database, as explained in Section~\ref{sec:problem}: the first copy opens a task for
editing, the second deletes it, and the first then saves its change.
% VERIFY in phase 3: every manual result in the tables is backed by a screenshot or was observed.

\subsection{The brief's test cases}

Table~\ref{tab:brief-tests} lists the seven test cases of the brief with their expected and actual results, the
steps by which each was carried out in the program, and the JUnit test that checks the same case at the service.

\begin{table*}[!tp]
  \caption{The brief's test cases}\label{tab:brief-tests}
  \centering\footnotesize
  \begin{tabularx}{\textwidth}{@{}L{0.025\textwidth}L{0.12\textwidth}YL{0.17\textwidth}L{0.2\textwidth}L{0.15\textwidth}@{}}
    \toprule
    No. & Test case & Steps in the program & Expected result & JUnit test & Actual result \\
    \midrule
    1 & Add a valid task
      & Click \emph{Add task} (or Ctrl+N), type a title and a description, and click \emph{Add task}.
      & The task is added and appears in the list, Pending.
      & \texttt{tc1AValidTaskIsAdded}
      & JUnit: passed. In the program: \gap{from Fig.~\ref{fig:app-03-task-added}} \\
    2 & Add a task with an empty title
      & Click \emph{Add task}, leave the title empty (or type only spaces), and click \emph{Add task}.
      & A validation message is shown, and no task is added.
      & \texttt{tc2ATaskWithAnEmptyTitle\-IsRefusedWithAMessage}
      & JUnit: passed. In the program: \gap{from Fig.~\ref{fig:app-04-empty-title}} \\
    3 & View all tasks
      & Add several tasks and give them all three statuses from their status menus.
      & All tasks are displayed, each with its status.
      & \texttt{tc3AllTasksAreListed}
      & JUnit: passed. In the program: \gap{from Fig.~\ref{fig:app-05-all-tasks}} \\
    4 & View an empty list
      & Start the program on a database with no tasks.
      & An appropriate message is shown in place of the list.
      & \texttt{tc4AnEmptyListIsEmpty}
      & JUnit: passed. In the program: \gap{from Fig.~\ref{fig:app-01-empty-list}} \\
    5 & Update an existing task
      & Click a task's \emph{Edit}, change its title, description and status, and click \emph{Save}.
      & The task's details are updated in the list.
      & \texttt{tc5AnExistingTaskIsUpdated}
      & JUnit: passed. In the program: \gap{from Fig.~\ref{fig:app-07-task-updated}} \\
    6 & Update a non-existing task
      & With two copies of the program open, click a task's \emph{Edit} in the first, delete the task in the second,
        then click \emph{Save} in the first.
      & Handled gracefully: a message says the task no longer exists, the list is refreshed, and no stack trace is
        shown.
      & \texttt{tc6UpdatingANonExisting\-TaskIsRefusedGracefully}
      & JUnit: passed. In the program: \gap{from Fig.~\ref{fig:app-08-update-missing}} \\
    7 & Delete an existing task
      & Click a task's \emph{Delete}, then \emph{Delete task} in the confirmation.
      & The task is removed from the list.
      & \texttt{tc7AnExistingTaskIsDeleted}
      & JUnit: passed. In the program: \gap{from Fig.~\ref{fig:app-10-task-deleted}} \\
    \bottomrule
  \end{tabularx}
\end{table*}

\subsection{Further test cases}

Table~\ref{tab:more-tests} lists the test cases we added for the requirements of Section~\ref{sec:problem} beyond the
brief's six, one row per JUnit test.

\begin{table*}[!tp]
  \caption{Further test cases}\label{tab:more-tests}
  \centering\footnotesize
  \begin{tabularx}{\textwidth}{@{}L{0.025\textwidth}YL{0.28\textwidth}L{0.27\textwidth}L{0.1\textwidth}@{}}
    \toprule
    No. & Test case & Expected result & JUnit test & Actual result \\
    \midrule
    \multicolumn{5}{@{}l}{\textit{The rules a task is saved under} (\texttt{TaskRulesTest})} \\
    8 & A title of spaces and tabs, or no title at all
      & Refused as an empty title; nothing saved.
      & \texttt{aBlankOrMissingTitleIsRefused} & Passed \\
    9 & Spaces and line breaks around the title and the description
      & Saved without them.
      & \texttt{theTitleAndDescriptionAreTrimmed} & Passed \\
    10 & No description
      & Kept as empty, never as null.
      & \texttt{noDescriptionIsKeptAsEmpty} & Passed \\
    11 & A title of 100 characters, then of 101
      & 100 accepted; 101 refused with a message giving the limit and the length.
      & \texttt{aTitleIsKeptWithinItsLimit} & Passed \\
    12 & A description of 1000 characters, then of 1001
      & 1000 accepted; 1001 refused with a message giving the limit and the length.
      & \texttt{aDescriptionIsKeptWithinItsLimit} & Passed \\
    13 & An update that breaks a rule
      & Refused, and the task is left as it was.
      & \texttt{aRefusedUpdateChangesNothing} & Passed \\
    14 & A status changed from the list, forward and back
      & The status changes; the title and description do not.
      & \texttt{aStatusIsChangedOnItsOwn} & Passed \\
    15 & A status change with no status
      & Refused with a message.
      & \texttt{aStatusChangeWithNoStatus\-IsRefused} & Passed \\
    16 & Changing the status of, deleting, or opening a task that no longer exists
      & Each refused with the message that the task no longer exists.
      & \texttt{aMissingTaskIsRefused\-ByEveryChange} & Passed \\
    17 & A change with no task chosen
      & Refused with ``Choose a task in the list first.''
      & \texttt{noTaskChosenIsRefused} & Passed \\
    \multicolumn{5}{@{}l}{\textit{Search, filter and counts} (\texttt{SearchAndCountsTest})} \\
    18 & Text that is in a title or a description, typed in another case
      & Every task holding it is found.
      & \texttt{searchFindsTheTextInATitle\-OrDescriptionInAnyCase} & Passed \\
    19 & No search text, or only spaces
      & Every task is listed.
      & \texttt{noTextMatchesEveryTask} & Passed \\
    20 & Text that is in no task
      & Nothing is found, and nothing fails.
      & \texttt{textInNoTaskFindsNothing} & Passed \\
    21 & Search text holding \texttt{\%} or \texttt{\_}
      & Matched as typed, not as wildcards.
      & \texttt{wildcardsAreMatchedAsTyped} & Passed \\
    22 & A status chosen in the filter
      & Only the tasks with that status are listed.
      & \texttt{aStatusFilterListsOnlyThatStatus} & Passed \\
    23 & Search text and a status together
      & Only the tasks matching both are listed.
      & \texttt{textAndAStatusNarrowThe\-ListTogether} & Passed \\
    24 & Tasks added, moved between statuses and deleted
      & Each count equals the number of tasks in its status after every change.
      & \texttt{theCountsFollowEveryChange} & Passed \\
    25 & No tasks
      & Every count is 0.
      & \texttt{anEmptyListCountsNothing} & Passed \\
    \multicolumn{5}{@{}l}{\textit{Two copies of the program} (\texttt{TwoCopiesTest})} \\
    26 & The other copy deletes a task; this copy then updates it, changes its status and deletes it
      & Each refused with the message that the task no longer exists; no error.
      & \texttt{aTaskDeletedInTheOtherCopy\-IsRefusedHereWithAMessage} & Passed \\
    27 & One copy adds and changes tasks
      & The other copy reads the changes at once.
      & \texttt{eachCopySeesTheOthersChanges} & Passed \\
    28 & The other copy deletes a task while this copy is saving a change to it
      & The delete waits for the save, then goes through; neither copy fails.
      & \texttt{aDeleteInTheOtherCopy\-WaitsWhileThisCopyIsSaving} & Passed \\
    \multicolumn{5}{@{}l}{\textit{Starting and restarting} (\texttt{RestartTest}, \texttt{FoundationTest})} \\
    29 & A task added and moved on, then the program closed and started again
      & The task is there as it was left.
      & \texttt{tasksSurviveARestart} & Passed \\
    30 & The program started on an empty database
      & Spring starts, Flyway makes the table, and Hibernate finds it as the entity expects.
      & \texttt{startsOnAFreshDatabase} & Passed \\
    \multicolumn{5}{@{}l}{\textit{The words the window shows, without a window} (\texttt{TaskTextTest})} \\
    31 & A time today, yesterday, earlier this year and in another year
      & Written as ``Today, 21:40'', ``Yesterday, 09:05'', ``16 Nov, 14:00'' and ``16 Nov 2025, 14:00''.
      & \texttt{aTimeIsWrittenTheWayThe\-ListWritesIt} & Passed \\
    32 & A moment just after midnight in Karachi, still the day before in UTC
      & Counted on the user's day, not on UTC's.
      & \texttt{theDayIsTheUsersNotUtcs} & Passed \\
    33 & A time inside a sentence
      & ``Today'' and ``Yesterday'' start in lower case; a date stays as it is.
      & \texttt{aTimeInsideASentence\-StartsInLowerCase} & Passed \\
    34 & The list's heading
      & ``All tasks'', or the status filtered by.
      & \texttt{theHeadingNamesWhatIsShown} & Passed \\
    35 & A number of tasks
      & Said with its status, and in the singular for one.
      & \texttt{aNumberOfTasksIsSaid\-WithItsStatus} & Passed \\
    36 & The line under the heading
      & Counts what the filter holds and, with a search, how many of those match.
      & \texttt{theSummarySaysWhat\-TheListShows} & Passed \\
    37 & An empty list
      & Says why: no tasks at all (test case 4), nothing matching the search, or nothing in the status.
      & \texttt{anEmptyListSaysWhy} & Passed \\
    38 & The character count in the form
      & Counted after trimming, as the service counts.
      & \texttt{charactersAreCounted\-AfterTrimming} & Passed \\
    \bottomrule
  \end{tabularx}
\end{table*}

\subsection{Results}

The last run of \texttt{./mvnw test} before this report was built ran all 38 tests, and every one passed:
\begin{lstlisting}[style=terminal]
[INFO] Tests run: 1, Failures: 0, Errors: 0, Skipped: 0, Time elapsed: 7.110 s -- in io.mcs.taskmanager.FoundationTest
[INFO] Tests run: 1, Failures: 0, Errors: 0, Skipped: 0, Time elapsed: 1.435 s -- in io.mcs.taskmanager.RestartTest
[INFO] Tests run: 7, Failures: 0, Errors: 0, Skipped: 0, Time elapsed: 0.798 s -- in io.mcs.taskmanager.services.BriefTestCasesTest
[INFO] Tests run: 8, Failures: 0, Errors: 0, Skipped: 0, Time elapsed: 0.752 s -- in io.mcs.taskmanager.services.SearchAndCountsTest
[INFO] Tests run: 10, Failures: 0, Errors: 0, Skipped: 0, Time elapsed: 0.528 s -- in io.mcs.taskmanager.services.TaskRulesTest
[INFO] Tests run: 3, Failures: 0, Errors: 0, Skipped: 0, Time elapsed: 2.234 s -- in io.mcs.taskmanager.TwoCopiesTest
[INFO] Tests run: 8, Failures: 0, Errors: 0, Skipped: 0, Time elapsed: 0.037 s -- in io.mcs.taskmanager.ui.TaskTextTest
[INFO]
[INFO] Results:
[INFO]
[INFO] Tests run: 38, Failures: 0, Errors: 0, Skipped: 0
[INFO]
[INFO] BUILD SUCCESS
[INFO] Total time:  16.443 s
\end{lstlisting}
These are the lines of that run that report the tests; Spring's log between them is left out. The whole run, Maven's
own start and the compilation included, took 16.4 seconds.

\screenshotwide{test-01-junit}
  {Running the tests with \texttt{./mvnw test}: the totals and BUILD SUCCESS.}
  {Terminal: ./mvnw test with the totals and BUILD SUCCESS}

Two faults were found and fixed while the program was being built, each together with the tests that now cover it:
the length limits that only Bean Validation enforced (test cases 11 and 12), and the save overtaken by a delete in
the other copy (test case 28), which that test reproduced before the fix. Both changes are shown in
Section~\ref{sec:git}.
